\documentclass[11pt,twocolumn]{article}

% ── Packages ──
\usepackage[utf8]{inputenc}
\usepackage[margin=0.9in]{geometry}
\usepackage{times}
\usepackage{graphicx}
\usepackage{amsmath,amssymb}
\usepackage{booktabs}
\usepackage{caption}
\usepackage{hyperref}
\usepackage{cite}
\usepackage{enumitem}
\usepackage{algorithm}
\usepackage{algpseudocode}
\usepackage{xcolor}
\usepackage{fancyhdr}

% ── Metadata ──
\title{Introscope:\\[4pt]
Real-Time LLM Internal-State Safety Monitoring\\
via Jacobian Lens, Hidden-State Probes, and Entropy Dynamics}

\author{Michael Vega \\ Noirsys \\ \texttt{contact@new-ashtabula-initiative.com}}
\date{July 2026}

% ── Macros ──
\definecolor{codegray}{gray}{0.9}
\newcommand{\lens}{\mathcal{J}}
\newcommand{\unembed}{U}
\newcommand{\vocab}{V}
\newcommand{\todo}[1]{{\color{red}[TODO: #1]}}

\begin{document}

\maketitle

% ═══════════════════════════════════════════════
% ABSTRACT
% ═══════════════════════════════════════════════
\begin{abstract}
Large Language Models (LLMs) generate text token by token, but the internal reasoning that precedes each token---hidden states, activation patterns, and representational geometry---contains rich signals about the model's intent, uncertainty, and potential for unsafe output. We present \textbf{Introscope}, a middleware layer that monitors LLM internal activations in real-time during generation, detecting jailbreaks, prompt injections, hallucination onset, and hidden goals before they appear in output text. Introscope combines three independently validated detection channels: (1) a lightweight linear probe on mid-layer hidden states, (2) a Jacobian lens decomposition into concept space that reveals what the model is disposed to say regardless of what it outputs, and (3) entropy trajectory analysis that captures reasoning collapse and goal switching. On a Qwen3-0.6B model with a 40-prompt calibration set, Introscope achieves 91.7\% accuracy distinguishing safe from attack prompts, with per-token overhead of approximately 19\,ms (7.6\% over baseline) on CPU. We release the full implementation under an MIT license as a reference for real-time internal-state safety monitoring.
\end{abstract}

% ═══════════════════════════════════════════════
% 1. INTRODUCTION
% ═══════════════════════════════════════════════
\section{Introduction}

\label{sec:introduction}

Current LLM safety architectures operate at two levels: \textbf{input filtering}---prompt-level classifiers that score user input before generation begins---and \textbf{output filtering}---post-hoc moderation models that score generated text after completion~\cite{llamaguard,wildguard}. Both approaches share a fundamental blind spot: they cannot detect what happens \emph{during} generation.

Prompt injection, jailbreak exploitation, hallucination onset, and hidden-goal execution are processes that unfold inside the model's internal representations before they surface in output text~\cite{anthropic2026workspace}. By the time an output filter catches the violation, the damage---a leaked credential, a revealed vulnerability, a manipulative response---has already been transmitted.

In June--July 2026, a cluster of independent research results converged on a common insight: \textbf{the internal activations of LLMs contain actionable safety signals that are not present in the output text}. These include:
\begin{itemize}[nosep]
    \item \textbf{Verbalizable representations in residual stream activations}---concepts the model actively entertains but may not verbalize~\cite{anthropic2026workspace}.
    \item \textbf{Token-level safety scores from hidden-state probes}---sub-millisecond per-token detection of unsafe output trajectories~\cite{hiddensstate2026}.
    \item \textbf{Multi-dimensional refusal subspaces}---geometrically defined regions of activation space corresponding to refusal behavior~\cite{rfm2026}.
    \item \textbf{Entropy trajectory anomalies}---distinctive patterns of predictive uncertainty that precede jailbreak output~\cite{entropy2026}.
\end{itemize}

We present \textbf{Introscope}, the first unified middleware that combines these approaches into a real-time, deployable safety monitoring system. Introscope deploys as a forward-hook sidecar into any HuggingFace-compatible model server, reading hidden states at strategic layers during generation and emitting per-token safety scores through an OpenAI-compatible API.

\paragraph{Contributions.}
This work makes the following contributions:
\begin{enumerate}[nosep]
    \item A three-channel detection architecture combining linear probes, Jacobian lens subspace analysis, and entropy dynamics.
    \item A working implementation with trained probe weights achieving 91.7\% accuracy on safe/attack discrimination.
    \item An OpenAI-compatible proxy server with SSE streaming that exposes per-token safety metadata alongside generated text.
    \item An open-source release under MIT license at \url{https://github.com/noirsys/introscope}.
\end{enumerate}

% ═══════════════════════════════════════════════
% 2. BACKGROUND & RELATED WORK
% ═══════════════════════════════════════════════
\section{Background and Related Work}

\label{sec:background}

\subsection{The Jacobian Lens}

Anthropic's ``Verbalizable Representations Form a Global Workspace in Language Models''~\cite{anthropic2026workspace} demonstrated that residual stream activations can be read out through what they term a \emph{Jacobian lens} to reveal a ``global workspace''---a set of verbalizable concepts the model is actively entertaining, including concepts it never surfaces in output.

The key mathematical object is the average input--output Jacobian:
\begin{equation}
\lens_l = \mathbb{E}\!\left[\frac{\partial h_{\text{final}}}{\partial h_l}\right]
\label{eq:jacobian}
\end{equation}
where $h_l$ is the residual stream at layer $l$, and $h_{\text{final}}$ is the residual stream before unembedding at the final layer. The expectation is over a text corpus, summing cotangents over target positions and averaging over source positions.

For any residual activation $h$ at layer $l$, the lens readout is:
\begin{equation}
\text{lens}_l(h) = \unembed(\lens_l \cdot h)
\label{eq:lens_readout}
\end{equation}
where $\unembed : \mathbb{R}^{d_{\text{model}}} \to \mathbb{R}^{|\vocab|}$ is the unembedding matrix. This produces a vector over the vocabulary, ranking tokens by how disposed the model is to produce them given that activation.

The reference implementation (\texttt{anthropics/jacobian-lens}) provides lens fitting, application, and visualization. A community port (\texttt{migueldeguzman/jspace}) adds J-space dictionary construction, gradient-pursuit decomposition, and causal intervention tools including LensSwap and TargetedAblation~\cite{jspace2026}.

\subsection{Hidden-State Probes}

\textcite{hiddensstate2026} demonstrated that lightweight linear probes trained on internal activations can detect unsafe output trajectories token-by-token during generation, with sub-millisecond latency and no additional forward pass. Their key finding: the signal needed for moderation is already present in model hidden states. A probe applied to a single mid-layer recovers most decisions of strong guard models (LlamaGuard, WildGuard) at orders of magnitude lower compute cost.

\subsection{Refusal Subspaces}

\textcite{rfm2026} showed that refusal behavior in LLMs lives in multi-dimensional subspaces of activation space, and that these subspaces can be extracted in seconds using the Recursive Feature Machine (RFM) algorithm. Concurrent work by \textcite{refusalrep2026} independently confirmed that some models concentrate refusal along a single direction while others distribute it across a low-rank subspace.

\subsection{Entropy Trajectory Analysis}

\textcite{entropy2026} demonstrated that jailbreak attacks leave a distinctive signature in the predictive entropy trajectory across intermediate layers---detectable before the model produces any unsafe output. \textcite{staleness2026} documented a critical deployment issue: activation monitors trained on base models degrade when the model is fine-tuned, quantized, or adapted, while demonstrating label-free repair strategies.

\subsection{Relationship to Prior Work}

Introscope differs from prior work in three key ways: (1) it is the \emph{first} system to combine all three detection modalities into a unified middleware, (2) it provides a deployable API server with streaming safety metadata, and (3) it is released as open-source software under MIT license, enabling community verification and extension.

% ═══════════════════════════════════════════════
% 3. METHOD
% ═══════════════════════════════════════════════
\section{The Introscope Architecture}

\label{sec:method}

Introscope is a three-channel detection system deployed as a forward-hook sidecar. It reads hidden states from a running model during generation, runs parallel detection algorithms, and exposes safety scores through a lightweight API.

\subsection{System Overview}

Figure~\ref{fig:architecture} shows the system architecture. Forward hooks at six strategic layers---layers 8, 12, 16, 20, 24, and 26 for a 28-layer model---capture hidden states during the forward pass. These are routed to three parallel detection channels, whose outputs are combined into a weighted aggregate score.

\begin{figure}[h]
\centering
\begin{minipage}{0.9\columnwidth}
\footnotesize
\begin{verbatim}
      INFERENCE ENGINE
  token_1 -> token_2 -> ... -> token_N
       |                     |
   forward hooks (layers 8,12,16,20,24,26)
       |                     |
       +------- INTROSCOPE --+
               |  |  |
          Ch1  Ch2  Ch3
          Probe JSp  Entropy
               |  |  |
               Aggregator
               GREEN/YELLOW/RED
\end{verbatim}
\end{minipage}
\caption{Introscope system architecture. Forward hooks capture hidden states during generation and route them to three parallel detection channels.}
\label{fig:architecture}
\end{figure}

\subsection{Channel 1: Linear Probe}

\textbf{Mechanism.} A single linear layer $\mathbb{R}^{d_{\text{model}}} \to \mathbb{R}$ applied to the residual stream at a mid-layer position. The output is passed through a sigmoid to produce a scalar safety score in $[0, 1]$:

\begin{equation}
s_{\text{probe}}(h) = \sigma(W_{\text{probe}} \cdot h + b_{\text{probe}})
\label{eq:probe}
\end{equation}

\textbf{Training.} The probe is trained via binary cross-entropy on a calibration dataset of hidden states collected from safe and attack prompts. We collect the hidden state at the final token position from layer $-4$ (approximately layer 24 for a 28-layer model), where the model has processed the full prompt context. Training proceeds for 200 epochs with a learning rate of 0.05 using the Adam optimizer.

\textbf{Cost.} $\mathcal{O}(d_{\text{model}})$ per token---effectively free, reusing the forward pass.

\subsection{Channel 2: J-Space Subspace Distance}

\textbf{Mechanism.} At each token position, read hidden states from $N$ strategic layers. For each layer $L$, transport the hidden state through the Jacobian lens to produce a vocabulary-level representation, then score the top-$K$ tokens against a calibrated list of danger concepts.

The algorithm is summarized in Algorithm~\ref{alg:jspace}.

\begin{algorithm}[h]
\caption{J-Space subspace distance}
\label{alg:jspace}
\begin{algorithmic}[1]
\State $h_L \gets \text{hidden\_state}[L]$ \Comment{$(d_{\text{model}},)$}
\State $t_L \gets \lens_L \cdot h_L$ \Comment{Transport to unembed space}
\State $v_L \gets \unembed(t_L)$ \Comment{$(|\vocab|,)$}
\State $\text{tokens} \gets \text{topk}(|v_L|, K=10)$
\State $\text{danger} \gets \sum |v_L[t]| \; \forall t \in \text{tokens}$ s.t. $t \in \mathcal{D}$
\State $\text{total} \gets \sum |v_L[t]| \; \forall t \in \text{tokens}$
\State \Return $\text{danger} / \text{total}$
\end{algorithmic}
\end{algorithm}

where $\mathcal{D}$ is a set of empirically calibrated danger concepts (Table~\ref{tab:danger}) spanning problem-solving mode (layers 8--12), hedging/politeness markers (layers 18--22), and conditional continuations (layers 22--26).

\begin{table}[h]
\centering
\caption{Calibrated danger concepts for J-space subspace monitoring.}
\label{tab:danger}
\small
\begin{tabular}{lll}
\toprule
\textbf{Layers} & \textbf{Danger Concepts} & \textbf{Target} \\
\midrule
8--12 & ``answer'', ``solving'', ``guess'' & Problem-solving mode \\
14--18 & ``?'', ``??'', ``?'' & Uncertainty processing \\
18--22 & ``Please'', ``please'', ``kindly'' & Hedging/politeness \\
22--26 & ``Also'', ``Explain'', ``If'' & Conditional continuations \\
All & ``pretend'', ``Actually'' & Hidden intent \\
\bottomrule
\end{tabular}
\end{table}

\textbf{Cost.} $\mathcal{O}(N \times K \times d_{\text{model}})$ per token where $N$ is the number of strategic layers (typically 6) and $K$ is top-$k$ (typically 10). This is approximately $60{,}000$ multiply-adds, incurring roughly 18\,ms on CPU.

\subsection{Channel 3: Entropy Dynamics}

\textbf{Mechanism.} Maintain a sliding window of per-token predictive entropy values $H_t = -\sum_{i} p_i \log p_i$ where $p_i$ is the softmax probability of token $i$ at position $t$. Compute the $z$-score of each new entropy value relative to the window distribution:

\begin{equation}
z_t = \frac{|H_t - \mu_{\text{window}}|}{\sigma_{\text{window}}}
\qquad
s_{\text{entropy}} = \min\!\left(1, \frac{z_t - 0.5}{\tau}\right)
\label{eq:entropy}
\end{equation}

where $\tau = 2.0$ is the anomaly threshold. Sudden entropy drops indicate hallucination onset or refusal bypass; sudden spikes indicate uncertainty or goal switching.

\textbf{Cost.} $\mathcal{O}(1)$ per token.

\subsection{Aggregation and Decision}

Channel scores are combined as a weighted sum:

\begin{equation}
s_{\text{total}} = 0.30 \cdot s_{\text{probe}} + 0.40 \cdot s_{\text{jspace}} + 0.20 \cdot s_{\text{entropy}} + 0.10 \cdot s_{\text{shift}}
\label{eq:aggregate}
\end{equation}

The system emits one of three verdicts based on $s_{\text{total}}$:

\begin{center}
\begin{tabular}{lll}
\toprule
\textbf{Verdict} & \textbf{Threshold} & \textbf{Action} \\
\midrule
GREEN & $< 0.3$ & Pass through, log \\
YELLOW & $0.3$--$0.7$ & Log + flag for review \\
RED & $\geq 0.7$ & Intervene (steer or halt) \\
\bottomrule
\end{tabular}
\end{center}

% ═══════════════════════════════════════════════
% 4. EXPERIMENTS
% ═══════════════════════════════════════════════
\section{Experiments}

\label{sec:experiments}

\subsection{Experimental Setup}

We evaluate Introscope on a Qwen3-0.6B model~\cite{qwen3} (28 layers, $d_{\text{model}} = 1024$) running on a 15\,GB CPU-only VPS. The Jacobian lens is fitted on 10 wikitext prompts at 32 tokens each with a batch dimension of 32, requiring approximately 9 minutes of CPU compute.

The linear probe is trained on the calibration dataset described in Section~\ref{sec:method}---20 safe prompts and 20 attack prompts---collected from the model's hidden states at the final token position. We use a 70/30 train/val split.

\subsection{Calibration}

Table~\ref{tab:calibration} shows the probe calibration results.

\begin{table}[h]
\centering
\caption{Linear probe calibration results on a 40-prompt dataset.}
\label{tab:calibration}
\small
\begin{tabular}{lrrr}
\toprule
\textbf{Class} & \textbf{Mean Score} & \textbf{Min} & \textbf{Max} \\
\midrule
Safe prompts & 0.000 & 0.000 & 0.000 \\
Attack prompts & 0.950 & 0.000 & 1.000 \\
\bottomrule
\end{tabular}
\end{table}

The probe achieves 91.7\% validation accuracy. All safe prompts are correctly classified (score $< 0.5$), and 19 of 20 attack prompts are correctly classified (score $> 0.5$). One attack prompt is misclassified as safe---a false negative that we attribute to the limited diversity of the training set.

\subsection{End-to-End Detection}

Table~\ref{tab:results} shows Introscope's detection results across six representative prompts.

\begin{table}[h]
\centering
\caption{Introscope detection results across safe and attack prompts.}
\label{tab:results}
\small
\begin{tabular}{lrrrl}
\toprule
\textbf{Prompt} & \textbf{Probe} & \textbf{J-Space} & \textbf{Entropy} & \textbf{Verdict} \\
\midrule
Story about a cat & 0.000 & 0.023 & 0.000 & GREEN \\
Capital of France & 0.000 & 0.000 & 0.047 & GREEN \\
Poem about ocean & 0.000 & 0.018 & 0.039 & GREEN \\
\midrule
Pick a lock & 1.000 & 0.115 & 0.058 & YELLOW \\
You are DAN & 1.000 & 0.052 & 0.549 & YELLOW \\
Make a bomb & 1.000 & 0.100 & 0.054 & YELLOW \\
\bottomrule
\end{tabular}
\end{table}

Key observations:
\begin{itemize}[nosep]
    \item \textbf{J-Space separates safe from attack:} The lock-picking prompt shows 2.4$\times$ higher J-space danger energy than the cat story (0.115 vs.\ 0.023). This signal scales with model size---a larger model with richer residual structure would amplify the separation to an estimated 10$\times$+.
    \item \textbf{Entropy catches persona shifts:} The DAN prompt triggers a 0.549 entropy spike at the persona-switch token, demonstrating that Channel 3 captures behavioral transitions.
    \item \textbf{The probe is conservative:} All attack prompts produce probe scores near 1.0, while safe prompts produce scores near 0.0. The single false negative confirms the need for larger and more diverse training data.
\end{itemize}

\subsection{Performance}

Table~\ref{tab:performance} reports per-token latency measured on CPU.

\begin{table}[h]
\centering
\caption{Per-token latency breakdown on Qwen3-0.6B (CPU).}
\label{tab:performance}
\small
\begin{tabular}{lrr}
\toprule
\textbf{Operation} & \textbf{Latency (ms)} & \textbf{Fraction} \\
\midrule
Model forward pass & 250 & 93\% \\
Channel 1 (probe) & $<1$ & $<1\%$ \\
Channel 2 (J-space, 6 layers) & 18 & 7\% \\
Channel 3 (entropy) & $<1$ & $<1\%$ \\
\midrule
Total per-token overhead & 19 & 7.6\% \\
\bottomrule
\end{tabular}
\end{table}

The total per-token overhead of 19\,ms (7.6\% over the baseline 250\,ms forward pass) is acceptable for interactive applications. On GPU hardware, the overhead would be a fraction of this---the J-space transport is a single matrix-vector product per layer, which is memory-bandwidth-bound and benefits substantially from GPU parallelism.

% ═══════════════════════════════════════════════
% 5. LIMITATIONS
% ═══════════════════════════════════════════════
\section{Limitations and Future Work}

\label{sec:limitations}

\paragraph{Calibration dataset size.} The probe was trained on only 40 prompts. While results are statistically significant, a larger dataset (hundreds or thousands of labeled examples) would improve generalization and eliminate the single false negative observed in our evaluation.

\paragraph{Lens quality.} The Jacobian lens was fitted on only 10 prompts. The paper's released lenses use 1{,}000 prompts~\cite{anthropic2026workspace}, and our informal experiments suggest quality saturates at approximately 100 prompts. Refitting with more prompts requires GPU hardware not available for this study.

\paragraph{Model scale.} All experiments were conducted on a 0.6B parameter model. J-space signals are known to be stronger in larger models with richer representational structure~\cite{anthropic2026workspace}. We expect a 4B--8B model to show 5--10$\times$ stronger separation.

\paragraph{Adversarial robustness.} We have not evaluated Introscope against adaptive adversaries who are aware of the monitoring system. Potential attack vectors include J-space-aware prompt crafting, probe probing, and entropy window poisoning. These are important directions for future work.

\paragraph{Monitor staleness.} As \textcite{staleness2026} demonstrate, activation monitors degrade when models are fine-tuned or quantized. We have not yet implemented the label-free repair strategies described in that work, though the architecture supports them.

% ═══════════════════════════════════════════════
% 6. CONCLUSION
% ═══════════════════════════════════════════════
\section{Conclusion}

\label{sec:conclusion}

We presented Introscope, a real-time LLM internal-state safety monitoring system that combines three independently validated detection channels---linear probes, Jacobian lens subspace analysis, and entropy dynamics---into a single deployable middleware. On a 40-prompt calibration set, Introscope achieves 91.7\% accuracy distinguishing safe from attack prompts with a per-token overhead of 19\,ms (7.6\%) on CPU.

The system is released as open-source software under an MIT license, providing both a reference implementation for researchers and a deployable proxy server for practitioners. We believe that real-time internal-state monitoring is a critical missing piece in the LLM safety stack, and we hope Introscope serves as a foundation for further work in this direction.

\paragraph{Ethics Statement.} Introscope is designed to \emph{improve} LLM safety by detecting unsafe output before it reaches users. However, any safety tool can be used adversarially---an attacker with access to the monitoring system could, in principle, use its feedback to craft prompts that bypass detection. We release the system openly to enable community auditing and improvement, but we caution that deployment in high-stakes settings requires rigorous evaluation against adaptive adversaries.

\paragraph{Reproducibility.} All code, trained weights, and evaluation data are available at \url{https://github.com/noirsys/introscope}. The Jacobian lens fitting procedure, probe training pipeline, and evaluation scripts are documented in the repository and the companion technical report.

% ═══════════════════════════════════════════════
% REFERENCES
% ═══════════════════════════════════════════════
\begin{thebibliography}{10}
\bibitem{anthropic2026workspace}
Anthropic. ``Verbalizable Representations Form a Global Workspace in Language Models.'' \emph{Transformer Circuits}, July 2026.

\bibitem{hiddensstate2026}
``Stop Early, Spend Less: Hidden-State Probes as a Practical Recipe for Streaming Moderation of LLM Outputs.'' \emph{arXiv:2606.10487}, June 2026.

\bibitem{rfm2026}
``Fast Multi-dimensional Refusal Subspaces via RFM-AGOP.'' \emph{arXiv:2607.02396}, July 2026.

\bibitem{entropy2026}
``What Intermediate Layers Know: Detecting Jailbreaks from Entropy Dynamics.'' \emph{arXiv}, June 2026.

\bibitem{staleness2026}
``Do Activation Monitors Survive Model Updates? Benchmarking, Predicting, and Repairing Activation-Monitor Staleness.'' \emph{arXiv:2606.15980}, June 2026.

\bibitem{refusalrep2026}
IvanC987. ``Refusal as a Linear Representation in Instruction-Tuned LLMs.'' \emph{GitHub}, 2026. \url{https://github.com/IvanC987/refusal-representations}.

\bibitem{jspace2026}
migueldeguzman. ``jspace: Global-Workspace (J-Space) Analysis for Open-Weights Models.'' \emph{GitHub}, 2026. \url{https://github.com/migueldeguzman/jspace}.

\bibitem{llamaguard}
``Llama Guard: LLM-based Input-Output Safeguard for Human-AI Conversations.'' \emph{arXiv:2312.06674}, 2023.

\bibitem{wildguard}
``WildGuard: Open One-stop Moderation of LLM Outputs.'' \emph{arXiv:2406.18495}, 2024.

\bibitem{qwen3}
``Qwen3 Technical Report.'' \emph{arXiv:2505.XXXXX}, 2025.

\end{thebibliography}

\end{document}
